\subsection{The exponent-sixteen fibres and the two-adic lifting argument}

\begin{statement}{E.14}{Proved}{Every CM matrix fibre survives the common exponent-sixteen refinement}
Let $J_{16,3}=(\gamma_3N\,N^{16})\cap(\gamma_3N\,N^3)$.
Its inverse image in $B_3$ is a finite-index normal subgroup
contained in $PB_3$. Every one of the 24 admissible CM matrix
pairs of E.1 occurs as a gentle shadow at this refinement,
with multiplier one modulo $576$.
\end{statement}
\begin{proof}
The exponent-sixteen collector has coordinates
$a\in(\Z/16)^{97}$ and $z\in(\Z/8)^{4656}$, with the product
in (E.13.2). The image of $N'N^6$, retaining $N^{\rm ab}/3$,
has normal coordinates
\[
 k=\prod_i n_i^{6u_i}\prod_{i<j}[n_i,n_j]^{3v_{ij}},
 \qquad (u,v)\in(\Z/8)^{4753},
\]
and exact product
\[
 (u,v)(u',v')=(u+u',\,v+v'-12(u_ju'_i)_{i<j}).
 \tag{E.14.1}
\]
The correction group is therefore nonabelian at this level;
the program retains its full product and inverse.

For each affine homology point of E.11 use the same ungauged
$S_0=\beta s\beta^{-1}$ and $R_0=S_0t$. With $\lambda=1+72h$,
the complete equations are
\[
 (S_0k)^2=1,\qquad
 (k^{-1}S_0t^{1+72h})^3=1,\qquad Hu=0\pmod8.
 \tag{E.14.2}
\]
Here $H$ is the two-row integer matrix of E.13.
The last condition is exact: the common kernel has
$F$-abelian image $48L_N$, and the correction has image
$6B_NHu$.

Write $x=(u,v,h)$ and let $\mathcal F(x)$ be the two
normal-coordinate order errors followed by $Hu$.
These are integer polynomials of degree at most two.
Their diagonal quadratic coefficients are divisible by six
and their mixed quadratic coefficients by twelve.
Indeed the sixth-power contribution of an automorphism
with abelian matrix $A$ and generator cochain $\tau$ is
\[
 2u_i\tau_i-(6u_i^2-u_i)A_{ai}A_{bi}
\]
in a central coordinate; reordering contributes multiples
of $12u_i u_j$. Cusp powers have the same divisibility,
since $12\binom h2=6h^2-6h$. Only linear homology
expressions are substituted into these formulas.
Thus $\mathcal F$ modulo four is exactly the affine system
of E.13.

Its homogeneous solutions are represented by
\[
 Vt+2Kw\pmod4,
 \qquad t\in\F_2^{801},\quad w\in\F_2^{827},
 \tag{E.14.3}
\]
where $K$ is the entire binary kernel and $V$ lifts the
kernel of the second-layer obstruction. Every column is
verified modulo four. The projection of $V\bmod2$ to
the 97 homology coordinates and the multiplier has rank
15. Select 15 columns spanning that projection and subtract
their combinations from the other columns. The remaining
786 columns have even homology and multiplier coordinates.
Twice a selected column is already represented by $2K$,
so this change preserves all solutions in (E.14.3).

Every direction with even homology and multiplier coordinates
changes $\mathcal F$ modulo eight linearly and independently
of the base point: the stated divisibilities kill its
quadratic cross terms. Hence all 786 remaining lower
parameters, all 827 upper parameters, and the arbitrary
correction $4y$ enter linearly; only the 15 selected bits
enter quadratically. The program
\texttt{de\_class2\_sixteen\_fibres.py} retains the full
linear system and evaluates all $32768$ choices of those
bits. In each of the 24 cases the linear correction rank
is 3953, and every constant, linear, and quadratic coefficient
vanishes modulo its image. All $32768$ choices survive.
A solution with $h=0$ is retained in every case and
substituted into both literal order equations and $Hu$.

For these two-primary coordinates and the three-primary
central correction $k_3$ of E.11, the common actual word is
\[
 k=\prod_i n_i^{6u_i}
       \prod_{i<j}[n_i,n_j]^{3v_{ij}+16(k_3)_{ij}}.
 \tag{E.14.4}
\]
The unimodular two-column calculation of E.11 supplies a
product $j$ of forty-eighth powers with
$\operatorname{ab}_F(kj)=0$. It is trivial in both factors.
Restoring the recorded cusp gauge gives
$f=sx^gS_0kx^{-g}j\in F'$.
The program \texttt{de\_fibre\_sixteen\_crt\_verify.py}
substitutes both full braid hexagons in both faithful
collectors, with 96 successful component checks.
Normal-block inversion retains the central term $-a_i a_j$.
Twelve additional literal conjugation checks include a
nontrivial change of section. Both Frattini ranks remain
97 because the corrections preserve the $G_6$ homology
point. This proves the full shadow and generation conditions.

The separate \texttt{de\_class2\_sixteen.py} finds no
lift of the 24 particular words retained in E.13. Its
projection fixes their upper corrections. Calculation
(E.14.3) permits all corrections over the homology point
and supplies the replacements above. These exact projection
maps explain the scopes of the two calculations.
\end{proof}

\begin{statement}{E.15}{Proved}{Every admissible CM pair survives the entire relative class-two tower}
For each of the 24 admissible CM matrix pairs there is a
compatible family of gentle shadows at the inverse images
in $B_3$ of
\[
             J_j=\gamma_3N\,N^j\qquad(6\mid j),
 \tag{E.15.1}
\]
ordered by divisibility, with multiplier $\lambda=1$.
All twelve extra matrix fibres therefore survive this
entire tower. The assertion specifies matrix fibres;
it does not assert extension of the particular normalized
words retained in E.13 or E.14.
\end{statement}
\begin{proof}
We construct the two-primary component first. Let
$\mathcal N_2$ be the free class-two pro-two group on
the 97 basis elements, in the extension of $Q$ obtained
by completing $N$ in this quotient. Let
$\mathcal U\subset\mathcal N_2$ have even abelian
coordinates. Multiplication by three is invertible over
$\Z_2$, so the elements $n_i^6,[n_i,n_j]^3$ give its
normal coordinates $(u,v)\in\Z_2^{4753}$, with product
(E.14.1). The polynomial bijection
\[
 Z=(u,z),\qquad z_{ij}=v_{ij}+6u_i u_j
\]
puts its operations in the form
\[
 Z*Z'=Z+Z'+\tfrac12[Z,Z'],\qquad
 [Z,Z']=(0,12u\wedge u').
 \tag{E.15.2}
\]
Every coefficient is integral. We use $\exp Z$ and
$\log$ for this explicit group-coordinate correspondence;
(E.15.2) justifies that notation at the prime two.

Conjugation by $S_0$ is an integral Lie automorphism $M_S$
with $M_S^2=1$. The first order equation is equivalent to
\[
                       G_S(Z)=(1+M_S)Z=0.
 \tag{E.15.3}
\]
In fact its logarithm is
$G_S(Z)+\frac12[G_S(Z),Z]$; the map
$Y\mapsto Y+\frac12[Y,Z]$ is an integral unipotent
automorphism. The target of $G_S$ is the saturated
submodule $\operatorname{Fix}(M_S)$.

Put $E=\log(R_0^3)$ and $r_*=\exp(-E/3)R_0$.
Here $R_0^3\in\mathcal U$, division by three is integral,
and $R_0$ commutes with its cube; thus $r_*^3=1$.
Write $M$ for its Lie action and
$P=(1+M+M^2)/3$ for the integral projection onto
$\operatorname{Fix}(M)$. For $k=\exp Z$, keep the
multiplier equal to one and write
\[
 R=k^{-1}S_0t=\exp(W)r_*,\qquad C=\log(R^3).
\]
The coordinates of $W$ and $C$ are integral polynomials
of degree at most two in $Z$, by (E.15.2).
The second order equation is equivalent to
\[
                              G_R(Z)=PC=0.
 \tag{E.15.4}
\]
To prove this integrally, the homology of $C$ is the
norm of the homology of $W$, so it is fixed by $M$ and
equals that of $PC$. Since $R^3$ commutes with $R$,
$(M-1)C=-[W,C]$. On $\ker P$ the inverse of $M-1$
is $-(M+2)/3$. With $Q_0=1-P$ this gives
\[
             C=PC+\frac{M+2}{3}Q_0[W,PC].
 \tag{E.15.5}
\]
The bracket uses only homology. Thus $PC=0$ is equivalent
to $C=0$, and (E.15.5) is an integral linear embedding
of the error $PC$, with left inverse $P$.

There is only one further independent commutator equation.
The integral matrix $H$ satisfies
\[
 H A_S=\begin{pmatrix}-2&-1\\3&2\end{pmatrix}H,
 \tag{E.15.6}
\]
where $A_S$ acts on $N^{\rm ab}$. Equivalently, apply
$B_N^{-1}$ to the interchange of the two coordinates
of $F^{\rm ab}$. The first row of $H(1+A_S)u$ is
$-(H_1+H_2)u$. Hence (E.15.3) and $H_2u=0$ imply
both commutator rows, and the converse holds. The exact
compressed system is the integral polynomial map
\[
 G:\Z_2^{4753}\longrightarrow
 \operatorname{Fix}(M_S)\oplus\operatorname{Fix}(M)\oplus\Z_2,
 \qquad G=(G_S,G_R,H_2u).
 \tag{E.15.7}
\]
The target lattices are primitively embedded. The
normal-coordinate changes and (E.15.5)--(E.15.6) recover
the full error vector from $G$, with primitive linear
part at a zero. At a point whose errors are divisible
by eight the same derivative comparison holds modulo
four, since the additional derivative terms multiply
those errors.

Here are the exact target and derivative ranks.
The cycle actions have traces $-7,-5,-5$ for $s,r,r^2$.
The exterior-square traces are $-24,15,15$, respectively,
from
$\operatorname{tr}(\Lambda^2A)=
((\operatorname{tr}A)^2-\operatorname{tr}(A^2))/2$.
Central cochains do not change a trace. It follows that
\[
 \operatorname{rank}\operatorname{Fix}(M_S)=2361,
 \quad \operatorname{rank}\operatorname{Fix}(M)=1591,
 \quad \operatorname{rank}\operatorname{target}(G)=3953.
 \tag{E.15.8}
\]
The program \texttt{de\_two\_hensel\_tangent.py}
checks these traces and (E.15.6). At each of the 24
points from E.14 it computes the full integral Jacobian
modulo four. The program \texttt{de\_two\_hensel\_fixed.py}
then deletes the multiplier column, recomputes the entire
binary kernel, and verifies all the ranks below for the
resulting 4753-variable system. Its binary kernel has
dimension 826. The divisibilities in the proof of E.14
give
\[
 \mathcal F(x+2e_i)-\mathcal F(x)=2J_i\pmod8
\]
and make $J\bmod4$ independent of $x$. The central
columns are the exterior-square norm matrices. All
homology columns are retained, with
five additional complete direction checks per point.
For the full binary kernel $K$, the exact ranks are
\[
 \operatorname{rank}_{\F_2}(J\bmod2)=3927,\qquad
 \operatorname{rank}_{\F_2}
 \left(J\bmod2\ \middle|\ \frac{JK}{2}\bmod2\right)=3953.
 \tag{E.15.9}
\]
The primitive comparison with (E.15.7) identifies these
with the ranks for $dG$. Its target has rank 3953, so
its elementary divisors consist of 3927 units and 26
elements of valuation one. Its image contains twice
the target lattice.

The lifting can now be given explicitly. At the integral
point $x_0$ retained in E.14, $G(x_0)$ is divisible
by eight. Choose a right inverse of $dG_{x_0}$ over
$\Q_2$ carrying the integral target into
$\frac12\Z_2^{4753}$; the elementary divisors give
such an inverse. The Newton correction belongs to
$4\Z_2^{4753}$. The new error is divisible by sixteen,
because $G$ is integral of degree at most two, and its
derivative is unchanged modulo four. Repeating, an
error of valuation $v\geq3$ is replaced by one of
valuation at least $2v-2$. The corrections converge to
an exact zero of $G$, with multiplier still equal to one.
Equations (E.15.3)--(E.15.6) give both exact orders
and zero $F$-abelianization for the correction $k$.
The retained cusp gauge commutes with $t$ and makes
the normalized $f$ have zero abelianization, by the
formula (E.15.12) below. The reduction on $N^{\rm ab}/2$ has
rank 97, so the pro-two Frattini argument gives an
automorphism at every finite level.

We give the gluing before applying the cusp gauge.
For the same CM matrix pair, take the ungauged
three-primary action of E.12, and take the identity
action at every other prime. In each component the
images of $s,r,t$ in $Q$ are $s,r,t$: the matrix
representatives $\beta,A$ centralize $s,r$ in $Q$,
and all their corrections belong to the normal kernel.
Thus these ungauged actions have the same identity
action on $Q$. The precise comparison of their target
groups is
\[
 \varprojlim_j B/(\gamma_3N\,N^j)
 \simeq
 \prod_{p,Q}\varprojlim_a B/(\gamma_3N\,N^{p^a}),
 \tag{E.15.10}
\]
where the product is the fibre product over $Q$.
On the normal kernel this is the Sylow decomposition
of the class-two completion of $N$: every finite
nilpotent quotient is the product of its Sylow groups.
For surjectivity at finitely many levels, lift their
common $Q$ coordinate and adjust independently in
these normal Sylow groups. Compactness gives (E.15.10).

We now specify one gauge in this common group. Let
$g_2\in\Z_2$ be the integral gauge retained in E.14.
For the limiting three-primary conjugator $\beta_3$,
let $g_3\in\Z_3$ be the difference of its second and
first $F$-abelian coordinates. The E.12 corrections
preserve these coordinates modulo three, so the
reduction of $g_3$ is that of the recorded integral
gauge. The image of $N$ in $F^{\rm ab}$ has basis matrix
$B_N=\left(\begin{smallmatrix}6&4\\0&2\end{smallmatrix}\right)$;
therefore every integral change in this difference is
even. The two-primary residue of the three-primary
gauge can accordingly be kept equal to that of $g_2$.
Choose $g\in\widehat\Z$ with these two primary
components and with $g_p=0$ for $p\ne2,3$.
The element $x=t^{-2}$ has order six in $Q$.
Consequently $x^g$ restricts to the specified cusp
gauge in the two- and three-primary extensions.
Apply its inner conjugation to the glued ungauged
action, obtaining $S',R'$ with $S'R'=t$.

Put $f=sS'$. Its common $Q$ coordinate is
\[
             q(f)=s x^{\bar g}s x^{-\bar g},
             \qquad \bar g=g\bmod6.
 \tag{E.15.11}
\]
In particular $f$ need not lie in the image of $N$.
The independent matrix program
\texttt{de\_adelic\_gauge.py} verifies the identity
ungauged $Q$ action and (E.15.11) for all 24 pairs.
For $\bar g=0,4,2$, respectively, its canonical matrix
representatives in $\mathrm{SL}_2(\Z/12)/\{\pm1\}$ are
\[
 I,\qquad
 \begin{pmatrix}1&4\\4&5\end{pmatrix},\qquad
 \begin{pmatrix}1&8\\8&5\end{pmatrix}.
\]
The exact abelianization formula is
\[
 \operatorname{ab}_F(sx^g Sx^{-g})
       =\operatorname{ab}_F(sS)+(-g,g).
 \tag{E.15.12}
\]
At two and three the selected gauges make this zero.
At every other prime the ungauged action is identity,
and $g_p=0$, so it is zero there also. In those
components the final action is inner by $x^g$; its
image can have nontrivial prime-to-$p$ order dividing
six, exactly as (E.15.11) requires. Hence $f$ has zero
full $F$-abelianization. Since $\gamma_3N\subset F'$,
the kernel of the resulting map to
$\widehat F^{\rm ab}=\widehat\Z^2$ is the image of
$\widehat F'$, and $f$ belongs to that image.
Both order equations hold in every component. The
exact hexagon equivalence of E.9 now supplies compatible
gentle shadows. Generation holds on each normal Sylow
kernel by its Frattini map at two and three and by
the specified inner maps at the other primes; it also
holds on $Q$. Each finite map is therefore an
automorphism. The two- and three-primary reductions
recover the prescribed CM matrix pair, proving the family.

Every subgroup in (E.15.1) contains $\gamma_3N$.
The map to $A_5$ in E.5 sends that subgroup onto
$A_5$, so its common refinement factors through none
of these quotients. This explicit comparison leaves
the cofinal lifting problem of E.2 unresolved.
\end{proof}

\begin{statement}{E.16}{Proved}{Simultaneous extension through every nonabelian simple quotient}
Let $\Sigma$ be any finite collection of epimorphisms
$\theta:B\twoheadrightarrow S_\theta$ to nonabelian finite
simple groups, and put
\[
 K_{j,\Sigma}=\gamma_3N\,N^j\cap
                \bigcap_{\theta\in\Sigma}\ker\theta,
 \qquad 6\mid j.
 \tag{E.16.1}
\]
Every one of the 24 CM matrix pairs has a compatible family
of gentle shadows at the inverse images of these subgroups
in $B_3$, with multiplier one, as $j$ increases by
divisibility and $\Sigma$ increases by inclusion.
Thus the class-two families of E.15 extend simultaneously
through all these nonabelian simple quotients.
For each fixed epimorphism $B\twoheadrightarrow R$ to
an arbitrary finite perfect group, they also extend
compatibly in $j$ through
$\gamma_3N\,N^j\cap\ker(B\to R)$.

The finite solvable quotient
$B/(\gamma_4N\,N^5)$ is not a quotient, through the natural
map from $B$, of any $B/K_{j,\Sigma}$. This gives a precise
remaining finite comparison outside this enlarged family.
\end{statement}
\begin{proof}
Write $E_j=B/(\gamma_3N\,N^j)$ and let $R_\Sigma$ be the
image of $B$ in $\prod_{\theta\in\Sigma}S_\theta$.
The latter is a product of nonabelian simple groups,
with repeated diagonal factors identified. To prove this,
adjoin the simple factors one at a time. The kernel of
projection to the previous factors, viewed in the new
simple group, is either the whole group or trivial.
In the first case the new group is a direct product.
In the second, the new factor is a simple quotient of
the preceding product and is the graph of an isomorphism
from one of its simple factors. This proves the claim
by induction. In particular $R_\Sigma$ is perfect.

The group $Q$ is solvable. Indeed,
$\mathrm{SL}_2(\Z/12)$ is the product of its reductions
modulo four and three. The kernel of reduction
$\mathrm{SL}_2(\Z/4)\to\mathrm{SL}_2(\F_2)$ is an
elementary abelian group of order eight, and the target
is $\mathfrak S_3$. The other factor has a normal
quaternion group of order eight with cyclic quotient
of order three. Passing to the quotient by $\{\pm I\}$
preserves solvability. Since the kernel of $E_j\to Q$
is nilpotent, $E_j$ is solvable as well.

The image of $B$ in $E_j\times R_\Sigma$ projects
onto both factors. For a subdirect product $H\subset U\times V$,
put $U_0=\{u:(u,1)\in H\}$ and
$V_0=\{v:(1,v)\in H\}$. Sending one coordinate of a
pair in $H$ to the other defines an isomorphism
$U/U_0\simeq V/V_0$. Here that common quotient is
both solvable and perfect, hence trivial: the derived
series of a solvable perfect group is simultaneously
constant and eventually trivial. Therefore
\[
                    B/K_{j,\Sigma}\simeq E_j\times R_\Sigma.
 \tag{E.16.2}
\]
The same argument applies to the image of $F'$.
First $F$ maps onto $R_\Sigma$, since $B/F\simeq
\mathfrak S_3$ and a quotient of a perfect group that
is solvable is trivial. Then $F'$ also maps onto
$R_\Sigma$, which is perfect. Its image in $E_j$ is
solvable, so its joint image is the product of its two
images.

For each E.15 shadow use its automorphism on $E_j$
and the identity on $R_\Sigma$. Under (E.16.2) the
conjugator is $(f_j,1)$, which is in the image of
$F'$ by the preceding paragraph. Both hexagons hold
coordinatewise with multiplier one, and the product
map is an automorphism. These constructions commute
with every transition map. Equivalently, compactness
applied to the joint images of $\widehat F'$ gives
one lift of the compatible pair $(\{f_j\},1)$.
This proves the claimed compatible family.
For one fixed arbitrary perfect target $R$, the same
subdirect-product and $F'$ arguments apply directly
to $E_j\times R$ and give the additional assertion.
For a collection of such targets the proof above
uses perfection of their \emph{joint image}; the
simple-factor induction proves exactly that property
for the simultaneous family in (E.16.1).

For the final comparison set $L=\gamma_4N\,N^5$.
It is characteristic in $N$, hence normal in $B$,
and has finite index because $N$ is finitely generated
and its class-three exponent-five quotient is finite.
The quotient $B/L$ is solvable. Every homomorphism
from (E.16.2) to a solvable group kills $R_\Sigma$,
since its image is perfect. Therefore a factorization
of the natural map $B\to B/L$ through (E.16.2)
would factor through $E_j$ and kill $\gamma_3N$.

Here is an explicit word that it does not kill.
Send the first two free generators of $N$ to
\[
 a=I+E_{12}+E_{34},\qquad b=I+E_{23}
                    \quad\hbox{in }U_4(\F_5),
\]
and send the others to the identity. The target has
class three and exponent five: commutators add
superdiagonal degrees, and $(I+A)^5=I$ for a strictly
upper triangular $4\times4$ matrix $A$ in
characteristic five. With $[a,b]=aba^{-1}b^{-1}$,
direct multiplication gives
$[[a,b],a]=I+2E_{14}\ne I$.
The program \texttt{de\_solvable\_barrier.py}
checks these literal matrices. Hence
$[[n_0,n_1],n_0]\in\gamma_3N\setminus L$,
contradicting that factorization. This proves the
claimed finite comparison and the exact scope of
the enlarged family.
\end{proof}

\subsection{A four-strand pentagon constraint on the CM fibres}

\begin{statement}{E.17}{Proved}{The complete Ferrari-homology pentagon test and its affine fibres}
Put $H_3=\F_3^2\rtimes\mathrm{SL}_2(\F_3)$, with the
standard left action on column vectors. The assignments
\[
 \begin{aligned}
 \rho_4(\sigma_1)&=(0,A),&
 A&=\begin{pmatrix}1&1\\0&1\end{pmatrix},\\
 \rho_4(\sigma_2)&=(0,B),&
 B&=\begin{pmatrix}1&0\\-1&1\end{pmatrix},\\
 \rho_4(\sigma_3)&=(e_1,A)
 \end{aligned}
 \tag{E.17.1}
\]
define an epimorphism $B_4\to H_3$.
Let $L_4=\ker\rho_4\cap PB_4$. Then $L_4$ is normal
in $B_4$, is contained in $PB_4$, and
$|B_4/L_4|=5184$.

The complete pentagon in $PB_4/L_4$ is a well-defined
condition on the $G_6$ conjugator coordinate used in E.1.
Among the 24 particular $G_6$ representatives retained
in E.14, exactly nine satisfy it. They have indices
\[
             3,4,5,7,8,14,15,17,23;
 \tag{E.17.2}
\]
the extra-coset indices among these are $4,7,15$.
The other fifteen particular representatives cannot
be projections of full $\widehat{\mathrm{GT}}$ elements.

Every one of the 24 \emph{whole CM matrix fibres}, however,
contains a gentle $G_6$ shadow with multiplier one
whose actual commutator word satisfies this pentagon.
In the 18-parameter three-primary presentation of E.10,
before quotienting by cusp conjugations, its equation
is one nonzero affine linear equation. Its solution
space has dimension 17 and contains points with full
generation rank 97 and with vanishing class-two cusp
obstructions of E.11. This is a finite pentagon test;
no lift through all four-strand quotients is asserted.
\end{statement}
\begin{proof}
We use the multiplication
$(u,M)(v,T)=(u+Mv,MT)$ in $H_3$.
The three assignments in (E.17.1) satisfy both adjacent
braid relations and the relation between the distant
generators, by direct multiplication. Their Ferrari
kernel elements
\[
 \xi_1=\sigma_3\sigma_1^{-1},\qquad
 \xi_2=\sigma_1\sigma_2\sigma_1
                  \sigma_3^{-1}\sigma_2^{-1}\sigma_1^{-1}
\]
map to $(e_1,I),(e_2,I)$, while $A,B$ generate
$\mathrm{SL}_2(\F_3)$. This proves surjectivity.
These are the actual free Ferrari generators and their
homology action: the source comparison is
\cite[(2.9) and (4.15)]{DE_NakamuraTsunogaiYasuda2010}.
The program \texttt{de\_ferrari\_pentagon.py} checks the
literal relations and enumerates the joint image with
the adjacent transpositions in $\mathfrak S_4$;
its order is 5184. This is precisely $B_4/L_4$.

Write $a_{ij}$ for the standard pure Artin generators.
We evaluate the complete group equation
\[
 \begin{aligned}
 &f(a_{12},a_{23}a_{24})f(a_{13}a_{23},a_{34})\\
 &\qquad=f(a_{23},a_{34})
       f(a_{12}a_{13},a_{24}a_{34})f(a_{12},a_{23}).
 \end{aligned}
 \tag{E.17.3}
\]
Each displayed pair defines an actual homomorphism
$F_2\to PB_4$, followed by $\rho_4$.
All five factors are pure, so the permutation component
of (E.17.3) is already the identity. We also retain
the evaluation at
$((\sigma_1\sigma_2)^3,\sigma_3^2)$, which is the
conjugator in the standard Drinfeld image of
$\sigma_3$.

The simultaneous images of the 97 literal generators
of $N$ under these six homomorphisms commute and have
exponent three. The program checks every generator
cube and all $\binom{97}{2}=4656$ commutators.
It follows that each homomorphism kills $N'N^6$.
Thus all six evaluations factor through the actual
quotient $F_2/(N'N^6)$, which is the pure free factor
in $G_6$. Corrections by the whole kernel have zero
image in these targets. Consequently a failure of
(E.17.3) at the computed representative is a failure
for every lift of that $G_6$ conjugator.
The retained E.14 word differs from the base word
evaluated here by elements of $N'N^6$, so this
comparison applies to the E.14 representatives
themselves. Complete multiplication of all five
factors gives exactly (E.17.2).

Here is the exhaustive affine reduction for whole
CM matrix fibres. Let $P=\binom{B_0}{A_0}$ be the
$194\times18$ parameter matrix retained in E.10.
A parameter $z\in\F_3^{18}$ changes the two
conjugators by $n(B_0z),n(A_0z)$. For an actual
integral change $b$ in the first conjugator, put
\[
 \ell(b)=\frac{
       \operatorname{ab}_{F,2}(n(b))-
       \operatorname{ab}_{F,1}(n(b))}{2}\pmod3.
\]
The numerator is even by the matrix $B_N$ of E.9.
The new gauge is $g+2\ell(b)$ modulo six.
If $f_\beta=s\beta s\beta^{-1}$ denotes the
ungauged base conjugator, the actual new word is
\[
 f_z=y^{g+2\ell(b)}
       (s n(b)s) f_\beta n(b)^{-1}
                  x^{-g-2\ell(b)}.
 \tag{E.17.4}
\]
This identity follows by inserting $s^2=1$ and
$sxs=y$; it is the precise gauge comparison used
in the calculation.

Embed each affine group in $3\times3$ matrices.
On the commuting exponent-three images of $N$,
the formulas
\[
 \log(I+D)=D+D^2,\qquad
 \exp(L)=I+L+2L^2\quad\text{over }\F_3
\]
are inverse and turn products into sums. Their
nilpotent matrices have cube zero, and every
commutator is zero. A commuting family of nilpotent
endomorphisms admits a common complete invariant flag:
choose a nonzero common kernel vector, pass to its
quotient, and repeat. In that basis these matrices
are strictly upper triangular, so every product of
three is zero. Expanding the two displayed polynomials
therefore proves $\log(uv)=\log(u)+\log(v)$ for
these commuting matrices, including characteristic three.
The joint linear map
recording the six logarithms of $n(b)$, those of
$s n(b)s$, and $\ell(b)$ has rank two on the
18 parameter columns. The program
\texttt{de\_ferrari\_fibres.py} retains that exact
matrix, its two-column section, and its
16-dimensional kernel. It enumerates all nine
values of the two-dimensional image. Three
independent literal-word substitutions per fibre
check (E.17.4) against the logarithmic evaluation.

For every one of the 24 fibres, the pentagon
residual on those nine values is an affine
linear map of rank one. Every value is retained
in the JSON certificate, and direct substitution
checks its equality with the retained linear
map. Exactly three values satisfy the pentagon.
Each also gives an invertible induced action
on Ferrari homology; its two columns are the
images of $\xi_1,\xi_2$ under the three
Drinfeld generator formulas. The program checks
both intertwining relations and all braid
relations after this substitution.

Finally, in the remaining 16 parameters the
program retains an actual point in every fibre
at which the five class-two cusp equations of
E.11 vanish and the complete $97\times97$
generation matrix is invertible. These are
positive witnesses, not exclusions from a
failed search. Their three-primary changes
are represented by four times the corresponding
integral $N$ words, so their two-primary
homology component is unchanged. The literal
cusp relation in $G_6$, both generation ranks,
the exact integral gauge, zero $F$-abelianization,
and all five pentagon factors are checked again.
The resulting commutator words and every
parameter are retained. This proves nonemptiness
in all 24 whole matrix fibres and the asserted
finite scope.
\end{proof}

\begin{statement}{E.18}{Proved}{The next Ferrari central layer adds no equation on these affine fibres}
Let $\mathcal H_3$ be the class-two exponent-three free group on two
generators, in coordinates $(u,v,z)\in\F_3^3$ with multiplication
\[
 (u,v,z)(u',v',z')=
 (u+u',v+v',z+z'+2(uv'-vu')).
\]
The standard action of $\mathrm{SL}_2(\F_3)$ on $(u,v)$ fixes $z$.
Put $\widetilde H_3=\mathcal H_3\rtimes\mathrm{SL}_2(\F_3)$.
The assignments
\[
 \sigma_1\longmapsto((0,0,0),A),\qquad
 \sigma_2\longmapsto((0,0,0),B),\qquad
 \sigma_3\longmapsto((1,0,2),A)
 \tag{E.18.1}
\]
define an epimorphism $\widetilde\rho_4:B_4\to\widetilde H_3$,
where $A,B$ are those in (E.17.1). The normal pure kernel
$\widetilde L_4=\ker\widetilde\rho_4\cap PB_4$ has index $15552$.
Reduction of the Heisenberg central coordinate gives the precise
quotient $\widetilde H_3\to H_3$ of E.17.

All six evaluations used in E.17 still factor through $G_6$.
On each of the 24 whole affine CM fibres, their joint parameter
image has dimension two, with exactly the same kernel as in E.17.
The complete pentagon has exactly the same zero set on this
image: three of its nine points. Every retained E.17 literal
witness satisfies the complete pentagon and all three braid
relations in $\widetilde H_3$, and its three braid images generate
that group. Thus this additional central layer preserves all
24 nonempty fibres and introduces no further affine equation.
\end{statement}
\begin{proof}
In the usual Baker--Campbell--Hausdorff coordinates of the free
class-two group, write $z_0$ for the central coordinate and
take the free generators to have $z_0=0$. The literal Ferrari
actions (4.15) of \cite{DE_NakamuraTsunogaiYasuda2010} have
central linear terms $(0,1)$ for $A$ and $(1,0)$ for $B$.
The exact change of coordinates
\[
                 z=z_0-u-2v
\]
removes both terms: for $\ell=(1,2)$, their rows are
$\ell A-\ell=(0,1)$ and $\ell B-\ell=(1,0)$.
Consequently the actual Ferrari generators become
$(e_1,2),(e_2,1)$, and their commutator is $(0,0,1)$.
This proves that (E.18.1) uses the actual Ferrari quotient.
Direct multiplication verifies the braid relations and the
four literal free-generator action formulas. These elements
generate the Heisenberg kernel and $\mathrm{SL}_2(\F_3)$,
proving surjectivity and the order $27\cdot24=648$.
The program \texttt{de\_ferrari\_class2\_model.py} also verifies
the Heisenberg multiplication and all 24 matrix actions
exhaustively, then enumerates the joint image with
$\mathfrak S_4$, of order $15552$.

For each of the six evaluation homomorphisms, the simultaneous
images of $N$ have order at most nine. The program
\texttt{de\_ferrari\_class2\_fibres.py} enumerates each such
subgroup and verifies that it is elementary abelian. In the
seven displayed coordinates $(a,b,c,d,u,v,z)$, subtracting
the identity gives additive coordinates on each of these
subgroups; every pair product is checked. All generator
cubes and all 4656 generator commutators are therefore
trivial. This proves the factorization through $G_6$.

Stack these additive coordinates for $n(B_0t)$ and
$s n(B_0t)s$, together with the gauge row $\ell(B_0t)$
of E.17. The retained exact matrix has rank two. Stacking
it with the old E.17 matrix still has rank two, so their
kernels in $\F_3^{18}$ coincide. Formula (E.17.4) remains
an exact literal-word identity in the larger group.
The calculation substitutes all nine values of the
two-dimensional image in every fibre. Its full
three-coordinate pentagon residual vanishes precisely
at the three points where the E.17 two-coordinate
residual vanishes. Three further literal substitutions
per fibre check the additive-coordinate evaluation.
Finally, each of the 24 saved E.17 words is substituted
directly; all five factors, all braid relations, and
generation of the order-648 group are checked. The
JSON retains every one of the $24\cdot9$ residuals,
and the compressed certificate retains the complete
parameter matrix, its section, and its kernel.
\end{proof}
\begin{statement}{E.19}{Proved}{Complete compatible corrections at the modulus-nine Ferrari test}
For $R=\mathbb Z/9$, let $\mathcal H(R)$ have coordinates
$(u,v,z)\in R^3$ and multiplication
\[
 (u,v,z)(u',v',z')=
 (u+u',v+v',z+z'+\tfrac12(uv'-vu')),
 \qquad \tfrac12=5\text{ in }R.
\]
Put $H_9=\mathcal H(R)\rtimes\mathrm{SL}_2(R)$, with
the linear action on $(u,v)$ and the identity on $z$.
The assignments (E.18.1), with the final central coordinate
changed from $2$ to $5$ and the matrices taken modulo nine,
define an epimorphism $\rho_{4,9}:B_4\to H_9$.
Its composition with reduction modulo three is precisely
$\widetilde\rho_4$ of E.18.

Let $\theta_i:F_2\to H_9$, $1\leq i\leq6$, be the
six actual evaluation maps of E.17. For
$B=C_2*C_3$, choose representatives $b_q$ of
$B/F_2\simeq\mathfrak S_3$ and set
\[
 K_6=N'N^6,\qquad
 K_{3,9}=K_6\cap
   \bigcap_{i=1}^6\bigcap_{q\in B/F_2}
       b_q^{-1}\ker(\theta_i)b_q,
 \qquad D_9=B/K_{3,9}.
 \tag{E.19.1}
\]
Thus $K_{3,9}$ is normal of finite index in $B$ and lies
in $F_2$. The kernel of $D_9\to G_6$ is
$\F_3^4$, and the image of $K_6\cap F_2'$ in that
kernel is $\F_3^3$.

For each of the 24 particular $G_6$ points retained in
E.17, all three commutator-correction parameters can be
solved simultaneously. The gentle norm equations have
rank two. The modulus-nine pentagon has rank three and
selects exactly one correction in this three-dimensional
image. The resulting literal words have zero integral
$F_2$-abelianization, satisfy both complete hexagons in
$D_9$, give surjective three-strand shadow maps, and
satisfy the complete pentagon in $H_9$. Their three
four-strand braid images satisfy all braid relations
and generate $H_9$.

This yields nonempty compatible finite tests for all
24 CM fibres. In this particular construction
$N'N^{18}\subset K_{3,9}$: the three-strand quotient
is an explicit quotient of $G_{18}$. The new condition
is the specified four-strand pentagon; no cofinal
extension follows from this statement.
\end{statement}
\begin{proof}
The exact coordinate change from the ordinary free
class-two coordinates is
$z=z_0+(u-v)/2$. It removes the central terms in the
two literal Ferrari actions over every ring in which
two is invertible. The free generators therefore map
to $(1,0,1/2)$ and $(0,1,-1/2)$. Direct multiplication
verifies their four action formulas and all three
braid relations over $R$. They generate
$\mathcal H(R)$, and the two elementary matrices
$A,B$ generate $\mathrm{SL}_2(R)$: row reduction
over the local ring $R$ first chooses a unit in the
first column, clears the other entry, and then
expresses the resulting diagonal matrix as a product
of elementary matrices. This proves the asserted
epimorphism. Its target has order
$9^3\cdot648=472392$.

Here is a faithful model for the core in (E.19.1).
An element $b\in B$ has permutation
$q\mapsto q\bar b$ of the six cosets and coordinates
\[
 h_{q,i}(b)=\theta_i(b_q b b_{q\bar b}^{-1}).
 \tag{E.19.2}
\]
The product rule is
$h_{q,i}(bc)=h_{q,i}(b)h_{q\bar b,i}(c)$.
Together with the original map to $G_6$, this is an
actual faithful representation of $D_9$; its kernel
is exactly (E.19.1). The program
\texttt{de\_ferrari\_mod9\_lift.py} constructs the
six representatives and checks $s^2=r^3=1$ in this
model. Its evaluation on $x,y$ is checked against
the six original maps $\theta_i$.

Reduction modulo three maps every $K_6$ coordinate
to the identity by E.18. The kernel of
$H_9\to\widetilde H_3$ is elementary abelian:
subtract the identity matrix and the zero
Heisenberg vector, divide all coordinates by three,
and reduce modulo three. Products add because
every cross term contains a factor nine. The
36 coordinates in (E.19.2) therefore give an
injective linear coordinate map on the image
of $K_6$.

The complete $K_6$ image is computed from the
4656 pairwise commutators of the 97 actual
Schreier generators $n_i$ of $N$, their sixth
powers, and their conjugates by the $n_i$.
These normally generate $N'N^6$ in $N$.
All pairwise commutator images are zero, and
all ninth powers are the identity. Thus this
particular core factors through $G_{18}$.
The span of the sixth-power images has rank
four; stability under every $n_i$ is checked.
This proves the exact kernel assertion.

The zero-abelianization restriction is imposed
integrally. With the previously retained matrix
$B_N=\left(\begin{smallmatrix}6&4\\0&2\end{smallmatrix}\right)$,
write
\[
       \operatorname{ab}_{F_2}(n_i)=B_N h_i,
       \qquad H=(h_i)_{0\leq i<97}.
\]
The columns numbered $0,8$ form a determinant-one
minor of $H$. They give an explicit integral
basis of the rank-95 lattice $\ker H$.
Modulo the commutator image, every element of
$K_6$ is represented by $\prod_i n_i^{6a_i}$.
Its $F_2$-abelianization vanishes exactly when
$a\in\ker H$. Evaluating this entire integral
lattice gives a rank-three image. The code
retains three literal words forming a basis;
their integral abelianizations are checked
to be zero. Their images are also checked
to be stable under $s,r$. This proves that
the allowed correction space is the entire
image of $K_6\cap F_2'$, with no omitted
power corrections.

For a retained E.17 conjugator $\beta$, put
\[
 S_0=\beta s\beta^{-1},\qquad
 t=sr,\qquad R_0=S_0t.
\]
Let $M=\im(K_6\cap F_2'\to D_9)$, and write
$S,R$ for conjugation by $S_0,R_0$ on $M$.
Replacing $S_0,R_0$ by $S_0k,k^{-1}R_0$ preserves
their product. Since $M$ is abelian, the two
torsion equations are exactly
\[
       (1+S)k=0,\qquad
       (1+R+R^2)k=R_0^3.
 \tag{E.19.3}
\]
All terms lie in the explicitly computed image
$M$; the code checks this by expressing them in
the retained basis. These equations have rank
two in every fibre.

Let $g$ be the exact integral cusp gauge of
the E.17 word. The corrected conjugator is
the literal commutator word
\[
          f=sx^gS_0k x^{-g}.
 \tag{E.19.4}
\]
For each evaluation map, changing $k$ changes
$f$ by an element of the elementary abelian
congruence kernel. Multiplying the five
pentagon factors therefore gives an affine
linear equation on $M$. The calculation
evaluates all $3^3=27$ correction values in
each of the 24 fibres, checking every
coefficient of this linear description.
The pentagon matrix has rank three; the
three braid-relation matrices have rank
zero. The joint system with (E.19.3) is
consistent of rank three, and so selects
one point of $M$ in every case.

The retained point is then verified by
complete literal substitution, without
using only its linear equations. The two
hexagons are checked in all 36 induced
coordinates, and hold in $G_6$ by E.17;
hence they hold in the faithful combined
quotient $D_9$. The word (E.19.4) has zero
integral abelianization because both its
base word and $k$ do.

For generation, the new three-strand map
is already onto $G_6$. Substitution of all
97 generators, their commutators and sixth
powers, followed by closure under their
conjugations, gives rank four in the kernel
of $D_9\to G_6$. Thus the map is onto $D_9$.
This supplies the finite gentle-shadow
generation condition, with multiplier one.
For the four-strand map, the first two
matrix images generate all 648 elements
of $\mathrm{SL}_2(R)$, as directly
enumerated. The two Ferrari generators
have trivial matrix part and determinant
nonzero modulo three on their two
Heisenberg homology coordinates. They
therefore generate the finite three-group
$\mathcal H(R)$ by its Frattini quotient.
The whole four-strand image is $H_9$.

Every correction word, source word, linear
matrix and right-hand side is retained
in the JSON and compressed certificates.
The distinction between these specified
finite tests and a compatible lift
through every four-strand quotient is
exactly the untested extension beyond
(E.19.1); it is not removed here.
\end{proof}


